\documentclass[12pt, a4paper]{article}
\usepackage[utf8]{inputenc}
\usepackage[T1]{fontenc}
\usepackage{amsmath, amssymb, amsthm}
\usepackage{mathrsfs}
\usepackage{hyperref}
\usepackage{geometry}
\usepackage{booktabs}
\usepackage{enumerate}
\geometry{margin=2.5cm}

\theoremstyle{plain}
\newtheorem{theorem}{Teorema}[section]
\newtheorem{lemma}[theorem]{Lema}
\newtheorem{proposition}[theorem]{Proposição}
\newtheorem{corollary}[theorem]{Corolário}
\theoremstyle{definition}
\newtheorem{definition}[theorem]{Definição}
\newtheorem{remark}[theorem]{Observação}
\theoremstyle{remark}
\newtheorem*{note}{Nota}

\title{Motor de Herança Estrutural:\\
	Formalização Algébrica da Fita-Dobra e do\\
	Mecanismo de Promoção de Pares de Goldbach}
\author{Tiago Bandeira\\
	\small\texttt{tiagobandeira.goldbach2025@gmail.com}}
\date{Maio de 2026}

\begin{document}
	
	\maketitle
	
	%--------------------------------------------------------------------
	\begin{note}[Nota do Autor~---~Distinção Central (rev.~2)]
		Esta revisão incorpora as definições consolidadas da
		\emph{Versão Definitiva} do Motor de Herança Estrutural e corrige
		as definições do motor presentes nas versões anteriores.
		A distinção central, que se propaga por \emph{todas} as definições,
		é:
		\medskip
		\begin{center}
			\begin{tabular}{lp{5.8cm}p{5cm}}
				\toprule
				Conceito & Definição & Quem garante Goldbach\\
				\midrule
				Par base $2M$ &
				Soma das \textbf{colunas} da fita:
				$a_k+b_k=2M$ &
				O sistema escolhe colunas com primos~---~Goldbach
				já é conhecido para $2M$.\\[4pt]
				Par alvo $2M{+}2$ &
				Soma das \textbf{diagonais} da fita:
				$a_{k+1}+b_k=2M{+}2$ &
				HR$^-$ certifica que $2M{+}2$ é Goldbach ao
				encontrar primos nos eixos de $W_i$.\\
				\bottomrule
			\end{tabular}
		\end{center}
	\end{note}
	
	%--------------------------------------------------------------------
	\begin{abstract}
		%--------------------------------------------------------------------
		Formalizamos o \textbf{Motor de Herança Estrutural}, um mecanismo
		dinâmico de cobertura dos números pares por somas de dois primos.
		O motor opera sobre dois objectos acoplados: a \textbf{fita-dobra}
		$\mathcal{F}_N^\varepsilon$ (uma matriz $2\times N$ de ímpares em
		percurso zigzag) e a grade $G_{3,C}$ (que fornece eventos restritivos
		via diagonais totalmente primas).
		
		A contribuição estrutural central deste artigo é a
		\textbf{separação em duas camadas}:
		\begin{itemize}
			\item \textbf{Camada algébrica} (incondicional): a bijeção de
			acoplamento $\Phi$, definida pela condição $3C=2N$, garante
			por álgebra pura a existência de pares simétricos deslocados
			em $W_i$ e a sua cobertura completa pela permutação
			deslizante $\sigma$.  Estes resultados independem
			inteiramente de primalidade.
			\item \textbf{Camada aritmética} (condicional): sob a
			\textbf{Hipótese Restritiva}~---~na forma fraca~(HR$^-$),
			que postula a existência de algum passo $t$ em que ambos os
			extremos de $W_i$ são primos, ou na forma
			forte~(HR$^+$), que exige uma tripla totalmente prima~---~o
			motor produz uma cadeia certificada de pares de Goldbach.
		\end{itemize}
		
		O \textbf{Teorema do Motor} estabelece, sob HR$^-$, que a cadeia
		$2M_0\to2M_0+2\to\cdots$ cobre todos os pares $\geq2M_0$ como somas
		de dois primos.  HR$^-$ e HR$^+$ são isoladas como as únicas
		hipóteses não provadas e serão o objecto do artigo seguinte, que as
		atacará via o peso oracle-free $\Omega_N^*$~\cite{peso} e a
		desigualdade estrutural~\cite{problemaB}.
	\end{abstract}
	
	\tableofcontents
	
	%====================================================================
	\section{Introdução}
	%====================================================================
	
	A Conjectura de Goldbach afirma que todo número par $>2$ é soma de
	dois primos.  Este trabalho integra-se na
	série~\cite{grade,invariante,peso,problemaB} e formaliza o
	\textbf{Motor de Herança Estrutural} introduzido em~\cite{motor},
	fornecendo as demonstrações algébricas que ali ficaram em aberto.
	
	A ideia central é substituir a pergunta estática
	\begin{center}
		\emph{``$2M$ é soma de dois primos?''}
	\end{center}
	pela pergunta dinâmica
	\begin{center}
		\emph{``dado que $2M$ é soma de dois primos (caso base),\\
			a estrutura garante que $2M+2$ também o é?''}
	\end{center}
	
	O motor opera sobre dois objectos acoplados:
	\begin{enumerate}
		\item A \textbf{fita-dobra} $\mathcal{F}_N^\varepsilon$, que
		organiza os ímpares em percurso zigzag e garante, por
		simetria posição-cega, a existência e a promoção de pares
		simétricos com qualquer soma prescrita.
		\item A \textbf{grade} $G_{3,C}$, que fornece \emph{eventos
			restritivos}~---~configurações totalmente primas numa janela
		$W_i$~---~que certificam a primalidade do par prometido.
	\end{enumerate}
	
	\medskip
	\noindent\textbf{Organização do artigo.}
	As Secções~\ref{sec:fita}--\ref{sec:sigma} constituem a
	\emph{camada algébrica}: todos os resultados são incondicionais e
	independem de primalidade.
	A Secção~\ref{sec:motor} introduz a \emph{camada aritmética}: as
	hipóteses restritivas HR$^-$/HR$^+$, o Motor de Herança e o Teorema
	principal.
	A Secção~\ref{sec:exemplos} ilustra ambas as camadas com exemplos
	numéricos explícitos.
	A Secção~\ref{sec:conclusao} sintetiza a separação e enuncia as
	lacunas em aberto.
	
	%====================================================================
	\section{A Fita-Dobra $\mathcal{F}_N^\varepsilon$}
	\label{sec:fita}
	%====================================================================
	
	\subsection{Definição}
	
	\begin{definition}[Fita-dobra]
		\label{def:fita}
		Fixados $N\geq1$ e $\varepsilon\in\{0,1\}$, a
		\textbf{fita-dobra} $\mathcal{F}_N^\varepsilon$ é a matriz
		$2\times N$ cujas células são:
		\[
		a_k = 2k-1,
		\qquad
		b_k = 4N - 2k + 1 - 2\varepsilon,
		\qquad k=1,\dots,N.
		\]
		A \textbf{linha~1} ($a_k$) percorre os ímpares da esquerda para a
		direita; a \textbf{linha~2} ($b_k$) percorre-os da direita para a
		esquerda (zigzag).  O parâmetro $\varepsilon=0$ é o
		\textbf{modo normal}; $\varepsilon=1$ é o \textbf{modo de
			repetição} (a última coluna da linha~2 repete o valor $b_N-2$,
		``tirando a vaga'' de um ímpar e reduzindo a soma em~$2$).
	\end{definition}
	
	\begin{remark}[Cardinalidade]
		Em ambos os modos, a fita contém exactamente $2N$ células.
		A repetição não altera a cardinalidade~---~substitui um valor por
		outro, mantendo o número de ``quadradinhos'' constante.
	\end{remark}
	
	\subsection{Soma constante (resultado incondicional)}
	
	\begin{proposition}[Soma constante]
		\label{prop:soma}
		Para todo $k\in\{1,\dots,N\}$ e todo $\varepsilon\in\{0,1\}$:
		\[
		a_k + b_k = 4N - 2\varepsilon =: 2M(\varepsilon).
		\]
		Toda coluna da fita soma o mesmo valor $2M(\varepsilon)$,
		independentemente da posição~$k$ e independentemente de
		primalidade.
	\end{proposition}
	
	\begin{proof}
		$a_k+b_k = (2k-1)+(4N-2k+1-2\varepsilon) = 4N-2\varepsilon$.
	\end{proof}
	
	\begin{remark}[Cegueira posicional]
		A Proposição~\ref{prop:soma} é puramente algébrica: a fita é
		\textbf{cega à posição} e \textbf{cega à primalidade}.  Não
		importa em que coluna~$k$ um par $(a_k,b_k)$ se encontra, nem se
		os seus elementos são ou não primos~---~a soma é sempre
		$2M(\varepsilon)$.
	\end{remark}
	
	\subsection{Dois pares consecutivos por fita}
	
	\begin{proposition}[Dois pares por fita]
		\label{prop:dois}
		A mesma fita física $\mathcal{F}_N$ representa simultaneamente
		dois números pares consecutivos:
		\[
		\mathcal{F}_N^0 \;\longrightarrow\; 2M^+ = 4N,
		\qquad
		\mathcal{F}_N^1 \;\longrightarrow\; 2M^- = 4N-2,
		\qquad
		2M^+ - 2M^- = 2.
		\]
	\end{proposition}
	
	\begin{proof}
		Directo da Proposição~\ref{prop:soma} com $\varepsilon=0$ e
		$\varepsilon=1$.
	\end{proof}
	
	\begin{remark}[Razão 2:1]
		Os ímpares estão distribuídos numa densidade que é o dobro da dos
		pares.  Por isso, cada coluna nova adicionada à fita cobre dois
		pares consecutivos, e a alternância de modo
		($\varepsilon=0\leftrightarrow1$) é o mecanismo que acessa cada
		um deles dentro da mesma estrutura física.
	\end{remark}
	
	\subsection{As diagonais da fita~---~Par alvo}
	\label{subsec:diagonais}
	
	Além dos pares colunares, a fita contém \textbf{pares diagonais}: cada
	par $(a_{k+1},\,b_k)$ cruza entre colunas adjacentes.
	
	\begin{proposition}[Diagonais da fita~---~par alvo]
		\label{prop:diagonal-alvo}
		Para todo $k\in\{1,\ldots,N-1\}$, na fita do par base
		($\varepsilon=1$, i.e.\ $2M^-=4N-2$):
		\[
		a_{k+1}+b_k^{(1)}
		=(2k+1)+(4N-2k-1)
		=4N=2M^+.
		\]
		Cada par diagonal $(a_{k+1},\,b_k^{(1)})$ soma exactamente
		$2M^+=2M^-+2$, independentemente de primalidade.
		O valor $2M^+$ é o \textbf{par alvo}~---~o número par
		imediatamente acima do par base $2M^-$.
		
		\medskip
		A fita representa assim dois papéis simultâneos:
		\begin{center}
			\begin{tabular}{lll}
				\toprule
				Objecto & Soma & Conceito\\
				\midrule
				Colunas: $a_k+b_k^{(1)}$ &
				$2M^-=4N-2$ & Par base (dado; Goldbach já conhecido)\\
				Diagonais: $a_{k+1}+b_k^{(1)}$ &
				$2M^+=4N$   & Par alvo (HR$^-$ certifica Goldbach)\\
				\bottomrule
			\end{tabular}
		\end{center}
	\end{proposition}
	
	\begin{proof}
		$a_{k+1}+b_k^{(1)}=(2(k+1)-1)+(4N-2k+1-2)=(2k+1)+(4N-2k-1)=4N$.
	\end{proof}
	
	\begin{remark}[Cobertura dos pares diagonais]
		Existem $N-1$ pares diagonais interiores na fita (o elemento $a_1=1$
		não forma par diagonal secundário válido, pois não existe $b_0$).
		Cada par diagonal é capturado pelo scanner e armazenado num eixo de
		$W_i$ (Secção~\ref{subsec:Wi-estrutura}).  A certificação de que
		$2M^+$ é Goldbach depende de encontrar uma diagonal com ambos os
		extremos primos~---~isto é precisamente HR$^-$.
	\end{remark}
	
	%====================================================================
	\section{Acoplamento $\Phi$ entre $G_{3,C}$ e $\mathcal{F}_N$}
	\label{sec:phi}
	%====================================================================
	
	\subsection{Condição de acoplamento}
	
	\begin{proposition}[Igualdade de cardinalidade]
		\label{prop:card}
		$G_{3,C}$ e $\mathcal{F}_N$ contêm o mesmo conjunto de ímpares
		se e somente se
		\[
		3C = 2N,
		\]
		ou seja, $N=\frac{3C}{2}$, o que exige $2\mid C$.
	\end{proposition}
	
	\begin{proof}
		$G_{3,C}$ tem $3C$ células e $\mathcal{F}_N$ tem $2N$ células.
		Para que ambas contenham exactamente os mesmos $3C=2N$ ímpares
		consecutivos é necessário e suficiente $3C=2N$.
	\end{proof}
	
	\subsection{Percurso zigzag}
	
	\begin{definition}[Percurso zigzag de $G_{3,C}$]
		\label{def:zigzag-grade}
		O \textbf{percurso zigzag} de $G_{3,C}$ é a sequência
		$\sigma_G:\{1,\ldots,3C\}\to\mathbb{N}_{\text{ímpar}}$ definida
		por:
		\[
		\sigma_G(k) = \begin{cases}
			2k-1
			& k=1,\ldots,C
			\quad\text{(linha~1, }\rightarrow\text{)}\\[4pt]
			2(2C-k)+1
			& k=C+1,\ldots,2C
			\quad\text{(linha~2, }\leftarrow\text{)}\\[4pt]
			2(k-2C)+4C-1
			& k=2C+1,\ldots,3C
			\quad\text{(linha~3, }\rightarrow\text{)}
		\end{cases}
		\]
	\end{definition}
	
	\begin{definition}[Percurso zigzag de $\mathcal{F}_N$]
		\label{def:zigzag-fita}
		O \textbf{percurso zigzag} de $\mathcal{F}_N$ é a sequência
		$\sigma_{\mathcal{F}}:\{1,\ldots,2N\}\to\mathbb{N}_{\text{ímpar}}$
		definida por:
		\[
		\sigma_{\mathcal{F}}(k) = \begin{cases}
			2k-1
			& k=1,\ldots,N
			\quad\text{(linha~1, }\rightarrow\text{)}\\[4pt]
			2(2N-k)+1
			& k=N+1,\ldots,2N
			\quad\text{(linha~2, }\leftarrow\text{)}
		\end{cases}
		\]
	\end{definition}
	
	\subsection{Bijeção de acoplamento (resultado incondicional)}
	
	\begin{definition}[Acoplamento $\Phi$]
		\label{def:phi}
		Dado $C\geq2$ com $2\mid C$, seja $N=\frac{3C}{2}$.
		O \textbf{acoplamento}
		\[
		\Phi: \{1,\ldots,N\} \longrightarrow G_{3,C}\times G_{3,C}
		\]
		associa a cada coluna $k$ da fita $\mathcal{F}_N$ o par de células
		correspondente em $G_{3,C}$, segundo a regra:
		\[
		\Phi(k) = \begin{cases}
			\bigl(f(1,\,k),\; f(3,\,C+1-k)\bigr)
			& k = 1,\ldots,C\\[6pt]
			\bigl(f(2,\,2C+1-k),\; f(2,\,k-C)\bigr)
			& k = C+1,\ldots,\tfrac{3C}{2}
		\end{cases}
		\]
		Para $k\leq C$, a fita emparelha a linha~1 da grade
		(percorrida $\rightarrow$) com a linha~3 (percorrida
		$\leftarrow$).  Para $k>C$, emparelha as extremidades da linha~2
		convergindo para o centro.
	\end{definition}
	
	\begin{proposition}[Acoplamento bem definido e soma constante]
		\label{prop:phi}
		Sob a condição $3C=2N$, $\Phi$ é uma bijeção bem definida.
		Além disso, para todo $k\in\{1,\ldots,N\}$, independentemente
		de primalidade:
		\[
		\Phi(k)_1 + \Phi(k)_2 = 4N = 2M^+.
		\]
	\end{proposition}
	
	\begin{proof}
		\textbf{Bijeção.}  Os pares $\Phi(k)$ para $k=1,\ldots,N$ cobrem
		exactamente os $2N=3C$ ímpares de $G_{3,C}$ sem repetição: para
		$k\leq C$ cobrem as linhas~1 e~3 ($2C$ elementos), e para $k>C$
		cobrem a linha~2 ($C$ elementos, emparelhados das extremidades
		para o centro).  Total: $2C+C=3C=2N$. \checkmark
		
		\textbf{Soma constante para $k\leq C$.}
		\[
		f(1,k)+f(3,C+1-k) = (2k-1)+(4C+2(C+1-k)-1)
		= (2k-1)+(6C-2k+1) = 6C = 4N.
		\]
		
		\textbf{Soma constante para $k=C+j$, $j=1,\ldots,C/2$.}
		Usando $f(2,c)=2(C+c)-1$:
		\[
		f(2,C+1-j)+f(2,j)
		= (4C-2j+1)+(2C+2j-1) = 6C = 4N. \quad\checkmark
		\]
	\end{proof}
	
	\subsection{Par deslocado garantido pelo acoplamento
		(resultado incondicional)}
	
	\begin{proposition}[Par deslocado garantido]
		\label{prop:par-deslocado}
		Seja $W_i$ qualquer janela preenchida pelo scanner na varredura
		de $G_{3,C}$, e seja $(e_1,e_3)$ os extremos de qualquer eixo
		$\mathcal{E}\in\{D_1,D_2,C_v,L_c\}$ de $W_i$.
		
		\medskip
		\textbf{Parte~1 --- Soma nos eixos (incondicional).}
		Via acoplamento $\Phi$ (Proposição~\ref{prop:phi}), cada par
		$(e_1,e_3)$ armazenado num eixo de $W_i$ satisfaz sempre~---~
		independentemente de primalidade~---:
		\[
		e_1 + e_3 = 4N = 2M^+.
		\]
		
		\medskip
		\textbf{Parte~2 --- Alinhamento após deslocamento (incondicional).}
		No próximo passo da fita ($\varepsilon:1\to0$ ou $N\to N+1$),
		a soma avança de $2$, independentemente de primalidade:
		o par prometido da iteração corrente torna-se par base para
		$2M^++2$, iniciando a próxima iteração do motor.
	\end{proposition}
	
	\begin{proof}
		\textbf{Parte~1.}  Pela Proposição~\ref{prop:phi}, para todo $k$,
		$\Phi(k)_1+\Phi(k)_2=4N=2M^+$.  O scanner preenche os eixos de
		$W_i$ exactamente com esses pares $\Phi(k)$, logo
		$e_1+e_3=2M^+$ em todo eixo. \qed
		
		\textbf{Parte~2.}  Por Definição~\ref{def:fita},
		$b_k=4N-2k+1-2\varepsilon$.  O deslocamento mínimo
		($\varepsilon\to0$ ou $N\to N+1$) incrementa $b_k\to b_k+2$ para
		todo~$k$ simultaneamente, logo $e_1+e_3\to e_1+e_3+2=2M^++2$. \qed
	\end{proof}
	
	\begin{remark}[A primalidade pertence à camada aritmética]
		A Proposição~\ref{prop:par-deslocado} garante a existência e o
		alinhamento do par $(e_1,e_3)$ para \emph{qualquer} eixo de
		\emph{qualquer} janela $W_i$, primo ou não.  A questão de
		\emph{quando} esse par é de primos~---~e de como isso certifica
		Goldbach~---~é tratada na Secção~\ref{sec:motor}, sob hipótese
		aritmética explícita.
	\end{remark}
	
	\subsection{A Janela $W_i$~---~Estrutura e Quatro Eixos de Simetria}
	\label{subsec:Wi-estrutura}
	
	\begin{definition}[Janela $W_i$~---~eixos de simetria]
		\label{def:Wi}
		A \textbf{janela} $W_i$ é uma submatriz $3\times3$ da grade
		$G_{3,C}$, centrada na coluna $i$, formada pelas três linhas e
		pelas colunas $i-1$, $i$, $i+1$.  Uma matriz $3\times3$ possui
		exactamente quatro eixos de simetria que passam pelo centro.
		Cada eixo conecta dois extremos opostos e define um
		\emph{slot} para um par diagonal da fita com soma $2M^+$ (par
		alvo):
		\[
		\begin{array}{|c|c|c|}
			\hline
			[0,0]\;D_1 & [0,1]\;C_v & [0,2]\;D_2\\
			\hline
			[1,0]\;L_c & [1,1]\;\text{PIVÔ} & [1,2]\;L_c\\
			\hline
			[2,0]\;D_2 & [2,1]\;C_v & [2,2]\;D_1\\
			\hline
		\end{array}
		\]
		\begin{center}
			\begin{tabular}{llc}
				\toprule
				Eixo & Posições & Soma nos extremos (par alvo)\\
				\midrule
				$D_1$ (diagonal principal)  &
				$[0,0]\leftrightarrow[2,2]$ & $E_1+E_2=2M^+$\\
				$D_2$ (diagonal secundária) &
				$[0,2]\leftrightarrow[2,0]$ & $E_1+E_2=2M^+$\\
				$C_v$ (coluna central)      &
				$[0,1]\leftrightarrow[2,1]$ & $E_1+E_2=2M^+$\\
				$L_c$ (linha central)       &
				$[1,0]\leftrightarrow[1,2]$ & $E_1+E_2=2M^+$\\
				\bottomrule
			\end{tabular}
		\end{center}
		Em cada eixo, o elemento menor ocupa a posição de entrada e o
		maior a posição oposta.  A célula central $[1,1]$ armazena o
		\textbf{pivô}~---~o elemento central da varredura do scanner.
		A soma $E_1+E_2=2M^+$ é garantida incondicionalmente pela
		Proposição~\ref{prop:diagonal-alvo}: é herança algébrica da
		construção da fita.
	\end{definition}
	
	\begin{corollary}[Soma constante nos eixos de $W_i$ --- corolário do acoplamento]
		\label{prop:soma-Wi}
		Sob a condição de acoplamento $3C=2N$, cada par $(E_1,E_2)$
		capturado pelo scanner e armazenado num eixo de $W_i$ satisfaz
		automaticamente
		\[
		E_1+E_2=4N=2M^+\quad\text{(par alvo),}
		\]
		independentemente de os elementos serem ou não primos.
		
		\medskip\noindent\textit{Demonstração.}
		O scanner preenche os eixos de $W_i$ com pares $\Phi(k)$ da
		bijeção de acoplamento.  Pela Proposição~\ref{prop:phi},
		$\Phi(k)_1+\Phi(k)_2=4N=2M^+$ para todo $k$, logo a soma
		é $2M^+$ em todo eixo. \qed
		
		\medskip\noindent
		Esta é a garantia algébrica incondicional do motor, e decorre
		directamente de $3C=2N$: é o acoplamento que impõe a soma
		constante, não uma propriedade independente da janela.
		A primalidade simultânea dos extremos é o único ingrediente
		aritmético adicional, capturado por HR$^-$.
	\end{corollary}
	
	\begin{remark}[Número de janelas por varredura]
		Como cada $W_i$ comporta até quatro diagonais (um por eixo) e
		existem $N-1$ pares diagonais interiores na fita, o número de
		janelas necessárias para cobrir todas as diagonais é
		$\lceil(N-1)/4\rceil$.  Quando $N-1$ não é divisível por~$4$,
		a janela final cobre o restante parcialmente, garantindo que
		nenhum par diagonal fique sem registo.
	\end{remark}
	
	%====================================================================
	\section{A Permutação Deslizante $\sigma$ e a sua Órbita}
	\label{sec:sigma}
	%====================================================================
	
	Esta secção é inteiramente incondicional: todos os resultados são
	álgebra pura, sem qualquer hipótese de primalidade.
	
	\subsection{Definição}
	
	\begin{definition}[Permutação deslizante $\sigma$]
		\label{def:sigma}
		Seja
		\[
		S = \bigl\{\sigma_G(2),\,\sigma_G(3),\,\ldots,\,
		\sigma_G(3C)\bigr\}
		\]
		o conjunto dos $3C-1$ elementos do percurso zigzag de $G_{3,C}$
		com excepção de $\sigma_G(1)=1$.  O elemento~$1$ é excluído
		porque não existe $p$ nos elementos da fita tal que $1+p=2M^+$:
		$1$ não forma nenhum par diagonal que soma o par alvo e,
		portanto, não pertence a nenhum eixo de $W_i$.
		
		A \textbf{permutação deslizante} $\sigma:S\to S$ é:
		\[
		\sigma\bigl(\sigma_G(k)\bigr) = \sigma_G(k-1),
		\qquad k = 2,\,3,\,\ldots,\,3C.
		\]
		Cada elemento avança uma posição no percurso zigzag, como uma
		peça de jogo que escorrega para o espaço vazio à esquerda.
	\end{definition}
	
	\begin{remark}[A permutação $\sigma$ como algoritmo de scanner com pivôs]
		\label{rem:scanner-pivos}
		A permutação deslizante $\sigma$ admite uma descrição algorítmica
		concreta~---~o \textbf{scanner com pivôs}~---~que é equivalente à
		Definição~\ref{def:sigma} e facilita a implementação computacional.
		
		\medskip
		\noindent\textbf{Algoritmo do Scanner:}
		\begin{enumerate}
			\item Posicionar o \textbf{pivô esquerdo} na segunda posição
			da lista de ímpares (pula $a_1=1$, pois~$1$ não forma par
			diagonal válido).
			\item Posicionar o \textbf{pivô direito} na extremidade final
			da lista.
			\item Registar o par $(E_{\text{esq}},E_{\text{dir}})$ num
			eixo de $W_i$.  O elemento central da lista é o pivô/sensor.
			\item Avançar o pivô esquerdo em $+1$ posição e recuar o
			pivô direito em $-1$ posição.
			\item Repetir até os pivôs se encontrarem.  Cada iteração
			preenche um eixo de $W_i$.
			\item Quando $W_i$ tiver os quatro eixos preenchidos,
			registá-la e iniciar nova $W_i$ para o bloco seguinte.
		\end{enumerate}
		
		\noindent\textbf{Equivalência com $\sigma$.}  Cada par de
		posições $(k_{\text{esq}},k_{\text{dir}})$ dos pivôs corresponde
		exactamente a um passo $t$ da órbita de $\sigma$ sobre
		$W_i$, com $(E_{\text{esq}},E_{\text{dir}})=(e_1(t),e_3(t))$
		e soma constante $e_1(t)+e_3(t)=2M^+$ (par alvo,
		Proposição~\ref{prop:soma-Wi}).  O pivô esquerdo pula~$1$
		porque~$1$ está excluído de~$S$ (Definição~\ref{def:sigma}).
	\end{remark}
	
	\begin{remark}[Por que excluir apenas o elemento~$1$]
		\label{rem:extremos}
		O único elemento excluído de~$S$ é $\sigma_G(1)=1$.  A razão
		é estritamente aritmética: $1$ não é primo e não existe $p$
		nos elementos da fita tal que $1+p=2M^+$, pelo que $1$ não
		pertence a nenhum par diagonal com soma $2M^+$ e não ocupa
		nenhum eixo de $W_i$.  O maior ímpar da grade~---~o último
		elemento de $S$~---~é o valor $\sigma_G(3C)$, que coincide
		com o maior ímpar da fita; quando $M$ é ímpar, a coluna
		central da fita repete o valor $a_N=b_N=M$ e o maior ímpar
		da grade é $2M^+-1-2=2M^--1$, não necessariamente $6C-1$.
		Em qualquer caso, esse elemento está incluído em~$S$ e o
		scanner posiciona o pivô direito sobre ele no primeiro passo
		(Passo~2 do Algoritmo do Scanner).
	\end{remark}
	
	\subsection{Cobertura completa dos pares diagonais
		(resultado incondicional)}
	
	\begin{proposition}[Órbita completa de $\sigma$]
		\label{prop:orbita}
		Seja $W_i$ a janela corrente do scanner na varredura de $G_{3,C}$.
		Para $t=1,\ldots,N-1$, independentemente de primalidade:
		\begin{enumerate}[(i)]
			\item \textbf{Ordenação.}  $\sigma^t$ preserva a ordem relativa
			dos elementos dentro de cada linha da grade.
			\item \textbf{Cobertura.}  A aplicação
			\[
			t \;\mapsto\;
			\text{par }(e_1(t),e_3(t))\text{ de }W_i
			\text{ sob }\sigma^t
			\]
			é uma bijecção de $\{1,\ldots,N-1\}$ sobre o conjunto
			dos $N-1$ pares diagonais interiores da fita
			$\mathcal{F}_N$.
			\item \textbf{Soma constante.}  Para todo $t$:
			\[
			e_1(t) + e_3(t) = 2M^+ = 4N.
			\]
		\end{enumerate}
	\end{proposition}
	
	\begin{proof}
		\textbf{(i)}  $\sigma$ é uma translação cíclica no índice $k$:
		$\sigma(\sigma_G(k))=\sigma_G(k-1)$.  Os elementos de cada linha
		são dispostos em ordem crescente (linhas~1 e~3) ou decrescente
		(linha~2) segundo~$k$; uma translação cíclica preserva essa ordem
		relativa.
		
		\textbf{(ii)}  O scanner inicia com o pivô esquerdo na posição~$2$
		e o pivô direito na posição~$3C$ do percurso zigzag
		(Passos~1--2 do Algoritmo do Scanner).  Após $t$ passos,
		os pivôs estão nas posições $t+1$ e $3C+1-t$.  À medida que
		$t$ percorre $\{1,\ldots,N-1\}$, os pares
		$(E_{\text{esq}},E_{\text{dir}})$ percorrem exactamente os
		$N-1$ pares diagonais interiores de $\mathcal{F}_N$,
		sem repetição.
		
		\textbf{(iii)}  Cada par $(e_1(t),e_3(t))$ é um par $\Phi(k)$
		da bijeção de acoplamento.  Pela Proposição~\ref{prop:phi}:
		$\Phi(k)_1+\Phi(k)_2=4N=2M^+$ para todo $k$.
		Logo $e_1(t)+e_3(t)=4N=2M^+$.
	\end{proof}
	
	\begin{corollary}[Um passo de $\sigma$ é um deslocamento mínimo]
		\label{cor:sigma-desloc}
		Cada passo de $\sigma$ corresponde a um deslocamento mínimo $+2$
		na fita.  A órbita completa de $\sigma$ percorre todos os $N-1$
		pares diagonais prometidos antes do próximo acoplamento.  Este
		percurso é completamente determinista e independe de primalidade.
	\end{corollary}
	
	\begin{remark}[Movimento dos pivôs e constância da soma]
		\label{rem:sigma-direcoes}
		A permutação $\sigma$ admite uma interpretação concreta em termos
		do scanner com pivôs: a cada passo~$t$, o pivô esquerdo avança
		uma posição na lista de ímpares e o pivô direito recua uma
		posição.  Como os ímpares estão espaçados de~$2$ em~$2$, isso
		traduz-se em:
		\[
		e_1(t{+}1)=e_1(t)+2,
		\qquad
		e_3(t{+}1)=e_3(t)-2.
		\]
		Os incrementos cancelam-se exactamente, garantindo a constância
		da soma a cada passo:
		\[
		e_1(t{+}1)+e_3(t{+}1)
		=\bigl(e_1(t)+2\bigr)+\bigl(e_3(t)-2\bigr)
		=e_1(t)+e_3(t)=2M^+.
		\]
		Esta é a razão pela qual cada par $(e_1(t),e_3(t))$ registado
		num eixo de $W_i$ soma sempre o par alvo $2M^+$,
		independentemente do passo~$t$: é uma consequência directa do
		movimento simétrico dos dois pivôs.
		
		\medskip
		\noindent\textbf{Configuração canónica e simulações.}
		Os elementos da grade não se movem durante a varredura do
		scanner: a grade está fixa na sua configuração canónica (ordenada)
		e é o par de pivôs que avança sobre ela, seleccionando a cada
		passo um novo par diagonal $(e_1(t),e_3(t))$ para um eixo de
		$W_i$.  O movimento dos elementos ocorre apenas nas simulações
		de configurações não canónicas~---~onde a lista de ímpares é
		embaralhada e reordenada gradualmente~---~que servem para
		estudar em que condições HR$^-$ ou HR$^+$ são activadas.
		
		\medskip
		\noindent\textbf{Nota sobre a natureza dos eventos.}
		Os passos $t=1,\ldots,N-1$ da órbita de $\sigma$ são
		\emph{eventos}~---~todas as configurações possíveis de $W_i$
		dentro de uma fita fixa $\mathcal{F}_N$.  Eles não são passos
		do motor: o motor avança com o crescimento da fita
		($N\to N{+}1$), e a grade é actualizada por propagação a cada
		novo passo.  A órbita de $\sigma$ enumera o
		\textbf{espaço de busca} da hipótese restritiva: quais os pares
		$(e_1,e_3)$ disponíveis para o evento~HR$^-$ numa fita fixa.
	\end{remark}
	
	\begin{proposition}[Pares diagonais da fita e unicidade do passo mínimo]
		\label{prop:pares-diagonais}
		Seja $\mathcal{F}_N$ a fita com $a_k=2k-1$ e
		$b_k^{(\varepsilon)}=4N-2k+1-2\varepsilon$.  Definem-se os
		\textbf{pares diagonais do scanner} de $\mathcal{F}_N^1$ como os
		pares capturados pelo scanner (Observação~\ref{rem:scanner-pivos}):
		\[
		(a_{k+1},\;b_k^{(1)}),\qquad k=1,\ldots,N-1.
		\]
		
		\medskip
		\noindent\textbf{Parte~1 --- Soma constante dos pares diagonais
			(incondicional).}
		\[
		a_{k+1}+b_k^{(1)}
		=(2k+1)+(4N-2k-1)
		=4N
		=2M^+.
		\]
		Cada par diagonal soma $2M^+$, independentemente de primalidade.
		
		\medskip
		\noindent\textbf{Parte~2 --- Cobertura bijectiva (incondicional).}
		A aplicação $t\mapsto(e_1(t),e_3(t))$ da
		Proposição~\ref{prop:orbita}(ii) é uma bijecção entre
		$\{1,\ldots,N-1\}$ e o conjunto dos $N-1$ pares diagonais
		interiores de $\mathcal{F}_N$.  Em particular, cada par diagonal
		da fita é \emph{apontado} por $W_i$ em exactamente um passo
		$t$ da órbita de $\sigma$.
		
		\medskip
		\noindent\textbf{Parte~3 --- Unicidade do passo mínimo $\delta=1$
			(incondicional).}
		Se $\sigma$ fosse substituída por $\sigma^\delta$ com $\delta>1$,
		ocorreriam duas quebras simultâneas:
		\begin{enumerate}[(i)]
			\item \emph{Cobertura parcial}: a órbita percorreria apenas
			$\lceil(N-1)/\delta\rceil$ dos $N-1$ pares diagonais,
			deixando $N-1-\lceil(N-1)/\delta\rceil>0$ pares sem
			cobertura~---~a bijecção da Parte~2 deixa de ser sobrejectiva.
			\item \emph{Dessincronização com o motor}: cada passo
			corresponderia a um salto de $\delta$ colunas na fita em vez
			de~$1$, pelo que o par prometido teria soma
			$2M^++2(\delta-1)>2M^+$ para $\delta>1$, quebrando a cadeia
			$2M^-\to2M^+$ do Teorema~\ref{thm:motor}.
		\end{enumerate}
		O passo $\delta=1$ é assim a \emph{única} escolha que mantém
		simultaneamente a cobertura bijectiva da Parte~2 e a
		sincronização $N\to N{+}1$ do motor.
	\end{proposition}
	
	\begin{proof}
		\textbf{Parte~1.}  Cálculo directo:
		$(2(k+1)-1)+(4N-2k+1-2)=(2k+1)+(4N-2k-1)=4N$.\checkmark
		
		\textbf{Parte~2.}  Os $N-1$ pares diagonais correspondem
		bijectivamente às colunas $k=1,\ldots,N-1$ da fita (a
		coluna~$N$ fornece o par simétrico central, não diagonal).  Pela
		Proposição~\ref{prop:orbita}(ii), a órbita de $\sigma$ cobre
		exactamente essas $N-1$ colunas sem repetição, logo a aplicação
		é uma bijecção.\checkmark
		
		\textbf{Parte~3(i).}  $\sigma^\delta$ tem órbita de comprimento
		$\lceil(N-1)/\delta\rceil$ sobre $W_i$, estritamente menor
		que $N-1$ para $\delta>1$.
		
		\textbf{Parte~3(ii).}  Por
		Observação~\ref{rem:sigma-direcoes}, cada passo de $\sigma$
		avança $e_1$ de~$2$ e recua $e_3$ de~$2$.  Após $\delta$ passos,
		$e_1$ avança $2\delta$ e $e_3$ recua $2\delta$.  O par prometido
		teria soma $4N+2(\delta-1)=2M^++2(\delta-1)>2M^+$ para $\delta>1$,
		implicando um passo do motor de $2\delta$ em vez de~$2$.\qed
	\end{proof}
	
	\begin{remark}[Unicidade do passo unitário]
		\label{rem:canonico}
		A Proposição~\ref{prop:pares-diagonais} Parte~3 mostra que
		o passo $\delta=1$ de $\sigma$ é a \emph{única} escolha que
		mantém simultaneamente a cobertura bijectiva completa dos $N-1$
		pares diagonais e a sincronização $2M^-\to2M^+$ do motor.
		Qualquer $\delta>1$ quebra as duas condições ao mesmo tempo.
	\end{remark}
	
	\subsection{Verificação numérica: $G_{3,10}$, $N=15$, $2M^-=58$}
	\label{subsec:verif-sigma}
	
	Para $C=10$, $N=15$, $2M^-=58$, $2M^+=60$.
	O scanner percorre os $N-1=14$ pares diagonais de $\mathcal{F}_{15}^1$.
	A tabela abaixo mostra os primeiros passos~---~cada linha é um passo
	$t$ do scanner, com o par $(e_1(t),\,e_3(t))$ registado num eixo de
	$W_i$.  A coluna ``par primo?'' é \emph{informação aritmética}~---~não
	interfere na álgebra das colunas anteriores.
	
	\[
	\begin{array}{c|cc|c|c}
		t & e_1(t) & e_3(t) & e_1+e_3 & \text{par primo?}\\
		\hline
		1 & 3  & 57 & 60 & \text{não (57 composto)}\\
		2 & 5  & 55 & 60 & \text{não (55 composto)}\\
		3 & 7  & 53 & 60 & \text{sim}\;\checkmark\\
		4 & 9  & 51 & 60 & \text{não}\\
		5 & 11 & 49 & 60 & \text{não (49 composto)}\\
		\vdots & \vdots & \vdots & \vdots & \vdots
	\end{array}
	\]
	
	A soma $e_1(t)+e_3(t)=60=2M^+$ é constante para todo~$t$
	(Corolário~\ref{prop:soma-Wi}).
	O par $(7,53)$ em $t=3$ activa HR$^-$: $7$ primo, $53$ primo,
	$7+53=60$~\checkmark.  O elemento~$1$ está excluído de~$S$ e não
	entra no scanner (pois não forma par diagonal com soma $2M^+$).
	O pivô direito inicia no maior ímpar da grade~---~neste exemplo,
	$\sigma_G(3C)=57$, que é o último elemento de LC~---~e está
	incluído em~$S$: a cobertura completa dos $N-1=14$ pares diagonais
	fecha correctamente após os pivôs se encontrarem.
	
	%====================================================================
	\section{O Mecanismo de Promoção}
	\label{sec:promocao}
	%====================================================================
	
	\begin{proposition}[Promoção por deslocamento]
		\label{prop:promocao}
		Seja $\mathcal{F}_N^\varepsilon$ a fita corrente com
		$2M(\varepsilon)=4N-2\varepsilon$, e $(e_1,e_3)$ qualquer par
		simétrico de $\mathcal{F}_N^\varepsilon$ com soma
		$e_1+e_3=2M(\varepsilon)$.  Um \textbf{deslocamento mínimo} é
		qualquer uma das operações:
		\begin{itemize}
			\item \textbf{Alternância de modo}:
			$\varepsilon\to1-\varepsilon$ com $N$ fixo;
			\item \textbf{Nova coluna}: $N\to N+1$ com $\varepsilon$ fixo.
		\end{itemize}
		
		\medskip
		\textbf{Parte~1 --- Fita (incondicional).}
		O deslocamento mínimo incrementa $b_k\to b_k+2$ para todo~$k$
		simultaneamente~---~independentemente de primalidade.  Logo, para
		qualquer par $(e_1,e_3)$, primo ou não:
		\[
		e_1 + e_3^{\text{novo}}
		= e_1 + (e_3+2)
		= 2M(\varepsilon)+2.
		\]
		
		\medskip
		\textbf{Parte~2 --- Grade (incondicional).}
		O deslocamento mínimo propaga-se para $G_{3,C}$ como uma
		\textbf{operação de pilha global}: fita e grade partilham os
		mesmos $3C=2N$ ímpares na mesma ordem, pelo que todo elemento de
		$G_{3,C}$ é reindexado simultaneamente.  Após o deslocamento,
		$f'(r,c) = f(r,c) + 2\delta(r)$, onde:
		\[
		\text{Alternância de modo:}\quad
		\delta(1)=0,\;\delta(2)=2,\;\delta(3)=0;
		\]
		\[
		\text{Nova coluna:}\quad
		\delta(1)=0,\;\delta(2)=0,\;\delta(3)=2.
		\]
		Em consequência, toda janela $W_i$ é reorganizada com novos
		valores~---~independentemente de primalidade.
	\end{proposition}
	
	\begin{proof}
		\textbf{Parte~1.}
		Por Definição~\ref{def:fita}, $b_k=4N-2k+1-2\varepsilon$.
		\begin{itemize}
			\item Alternância $\varepsilon\to0$ (se $\varepsilon=1$):
			$b_k^{\text{novo}}=4N-2k+1=b_k+2$. \checkmark
			\item Nova coluna $N\to N+1$ com $\varepsilon\to1$:
			$b_k^{\text{novo}}=4(N+1)-2k+1-2=b_k+4-2=b_k+2$. \checkmark
		\end{itemize}
		
		\textbf{Parte~2.}
		Os $3C=2N$ ímpares de $G_{3,C}$ e de $\mathcal{F}_N$ são os
		mesmos, dispostos em percursos zigzag compatíveis (Secção~\ref{sec:phi}).
		O deslocamento introduz um novo ímpar no extremo superior
		(ou repete um), empurrando todos os outros uma posição na ordem.
		Como a grade preenche as células linha a linha em ordem crescente,
		o empurrão global redistribui os valores por todas as linhas e
		colunas, reorganizando toda janela $W_i$. \qed
	\end{proof}
	
	\begin{remark}[O mecanismo é uma ferramenta algébrica]
		A Proposição~\ref{prop:promocao} mostra que a fita e a grade
		têm um mecanismo de deslocamento completamente determinista.
		Esse mecanismo funciona para \emph{qualquer} par $(e_1,e_3)$,
		seja ele de primos ou não.  O motor de herança usa esta ferramenta
		sob uma hipótese aritmética adicional (Secção~\ref{sec:motor}):
		a de que existe algum passo $t$ em que o par é efectivamente de
		primos.
	\end{remark}
	
	%====================================================================
	\section{Par Base e Par Prometido}
	\label{sec:pares}
	%====================================================================
	
	Esta secção introduz as definições que ligam a camada algébrica ao
	motor de herança.  O par base e o par prometido são definidos em
	termos de primalidade, mas a sua existência como pares simétricos
	deslocados é garantida incondicionalmente pelas secções anteriores.
	
	\begin{definition}[Par base]
		\label{def:base}
		Um par $(a_s,b_s)$ de $\mathcal{F}_N^1$ é um \textbf{par base}
		para $2M^-$ se $a_s$ e $b_s$ são ambos primos.  A escolha da
		coluna $s$ é livre~---~independente de qualquer evento restritivo.
		A existência de um par base é a afirmação de que $2M^-$ é soma
		de dois primos (caso base de Goldbach); o par base não é
		construído pelo motor, mas é o ponto de partida a partir do qual
		o motor opera.
	\end{definition}
	
	\begin{definition}[Par prometido]
		\label{def:prometido}
		O \textbf{par prometido} para $2M^+$ é um par $(p_1,p_3)$ de
		$\mathcal{F}_N^0$ tal que:
		\begin{enumerate}[(i)]
			\item $p_1+p_3=2M^+=4N$
			(garantido pela álgebra do acoplamento $\Phi$~---~incondicional);
			\item $(p_1,p_3)$ são os extremos de $W_i$ para algum
			passo $t$ da permutação deslizante $\sigma$
			(a associação com $W_i$ é incondicional~---~
			Proposição~\ref{prop:orbita});
			\item $p_1$ e $p_3$ são ambos primos
			(esta é a condição aritmética, postulada pela Hipótese
			Restritiva~HR$^-$~---~condicional).
		\end{enumerate}
		Quando existe, o par prometido é o certificado de Goldbach para
		$2M^+$: após o deslocamento mínimo da fita (passo $+2$), ele
		torna-se o par base para $2M^++2$.
	\end{definition}
	
	\begin{remark}[Separação das condições]
		As condições~(i) e~(ii) da Definição~\ref{def:prometido} são
		incondicionais: qualquer par de extremos de $W_i$ sob
		$\sigma^t$ satisfaz automaticamente $p_1+p_3=2M^+$ e está
		associado a $W_i$ via a órbita de $\sigma$.  Só a
		condição~(iii)~---~a primalidade simultânea de $p_1$ e $p_3$~---~é
		hipotética.  Esta separação é a essência da arquitectura do motor.
	\end{remark}
	
	\begin{proposition}[Independência entre par base e par prometido]
		\label{prop:indep}
		O par base $(a_s,b_s)\in\mathcal{F}_N^1$ e o par prometido
		$(p_1,p_3)\in\mathcal{F}_N^0$ são \textbf{independentes}: cada
		um soma o respectivo $2M(\varepsilon)$ por
		Proposição~\ref{prop:soma}, mas são determinados por mecanismos
		distintos~---~o par base por qualquer coluna prima da fita, o par
		prometido pelo evento restritivo de $W_i$.  Em particular,
		não precisam partilhar elementos.
	\end{proposition}
	
	\begin{proof}
		O par base existe em $\mathcal{F}_N^1$ ($\varepsilon=1$) e o par
		prometido em $\mathcal{F}_N^0$ ($\varepsilon=0$); os dois modos
		são distintos e independentes por construção.
	\end{proof}
	
	%====================================================================
	\section{O Motor de Herança: Hipóteses e Teorema}
	\label{sec:motor}
	%====================================================================
	
	As Secções~\ref{sec:fita}--\ref{sec:pares} estabeleceram a camada
	algébrica: toda a geometria dos pares deslocados, a cobertura pela
	órbita de $\sigma$, e o mecanismo de promoção são incondicionais.
	Esta secção introduz a camada aritmética: as hipóteses que postulam
	a existência de pares de primos, e o teorema que delas decorre.
	
	\subsection{As duas hipóteses restritivas}
	
	\begin{definition}[Hipóteses Restritivas]
		\label{def:HR}
		Fixado $C$, seja $(e_1(t),\,e_3(t))$ o $t$-ésimo par diagonal
		capturado pelo scanner na varredura de $G_{3,C}$, armazenado
		num eixo de alguma janela $W_i$.
		
		\medskip
		\textbf{Hipótese Restritiva Fraca~(HR$^-$).}
		\[
		\exists\, t\in\{1,\ldots,N-1\}:
		\quad e_1(t)\in\mathbb{P}
		\;\;\text{e}\;\;
		e_3(t)\in\mathbb{P}.
		\]
		Em algum passo da varredura, \emph{ambos} os extremos de um
		eixo de $W_i$ são primos.  Esta é a condição mínima para o motor:
		o par $(e_1(t),e_3(t))$ é então um par prometido
		(Definição~\ref{def:prometido}) que, após o deslocamento mínimo,
		torna-se par base para $2M^-+2=2M^+$.
		
		\medskip
		\textbf{Hipótese Restritiva Forte~(HR$^+$).}
		\[
		\exists\, t\in\{1,\ldots,N-1\}:
		\quad \bigl(e_1(t),\,p_2(t),\,e_3(t)\bigr)
		\text{ é totalmente prima.}
		\]
		Existe um passo em que a \emph{tripla completa} de um eixo de
		$W_i$ (diagonal $D_1$, $D_2$ ou coluna central $C_v$) é
		inteiramente formada por primos.  HR$^+$ implica HR$^-$.
		Esta condição mais forte conecta ao peso oracle-free
		$\Omega_N^*$~\cite{peso} e à desigualdade estrutural
		$\mu_{\text{comp}}>\mu_{\text{global}}>\mu_{\text{primo}}$~\cite{problemaB}.
	\end{definition}
	
	\begin{remark}[Relação entre HR$^-$ e HR$^+$]
		HR$^+$ implica HR$^-$ (uma tripla totalmente prima fornece em
		particular dois extremos primos).  O recíproco é falso em geral.
		O motor opera sob HR$^-$; HR$^+$ é a ferramenta para
		demonstrações mais profundas~---~nomeadamente a prova da própria
		HR$^-$ via densidade de eventos restritivos e a recorrência da
		cadeia de Markov associada.
	\end{remark}
	
	\subsection{Teorema do Motor de Herança}
	
	\begin{theorem}[Motor de Herança Estrutural]
		\label{thm:motor}
		\textit{Hipóteses:}
		\begin{itemize}
			\item[\textbf{(HI)}] Existe $2M_0$ com par base
			$(a_s,b_s)\in\mathcal{F}_{N_0}^1$ (caso base de Goldbach)
			e evento restritivo inicial em $G_{3,C_0}$.
			\item[\textbf{(HR$^-$)}] Para toda iteração do motor, existe
			pelo menos um passo $t\in\{1,\ldots,N-1\}$ tal que
			$(e_1(t),e_3(t))$ de $W_i$ são ambos primos.
		\end{itemize}
		
		\textit{Conclusão.}  A cadeia
		\[
		2M_0 \;\longrightarrow\; 2M_0+2 \;\longrightarrow\; 2M_0+4
		\;\longrightarrow\; \cdots
		\]
		cobre todos os pares $\geq 2M_0$ como somas de dois primos.
	\end{theorem}
	
	\begin{proof}
		Procedemos por indução, usando a camada algébrica das
		Secções~\ref{sec:fita}--\ref{sec:pares} em cada passo.
		
		\textbf{Base.}  Por~HI, $2M_0$ é Goldbach via par base
		$(a_s,b_s)\in\mathcal{F}_{N_0}^1$.  Por~HR$^-$, existe $t_0$
		tal que $(e_1(t_0),e_3(t_0))$ são ambos primos; pela
		Definição~\ref{def:prometido}, este par é o par prometido
		$(p_1,p_3)$ em $\mathcal{F}_{N_0}^0$ com
		$p_1+p_3=2M_0+2$.
		
		\textbf{Promoção.}  Na iteração seguinte $\mathcal{F}_{N_1}^1$,
		o par $(p_1,p_3)$ é promovido a par base para $2M_0+2$
		pela Proposição~\ref{prop:promocao} Parte~1~---~certificando
		Goldbach para $2M_0+2$.
		
		\textbf{Regeneração.}  Por~HR$^-$ aplicada à nova configuração,
		existe $t_1$ tal que $(e_1(t_1),e_3(t_1))$ são ambos primos;
		este par é o novo par prometido $(q_1,q_3)$ com
		$q_1+q_3=2M_0+4$.
		
		\textbf{Indução.}  Repetindo promoção e regeneração, cada par
		$2M_0+2n$ é certificado como soma de dois primos.
	\end{proof}
	
	\begin{remark}[Estatuto de cada componente]
		\label{rem:estatuto}
		\begin{center}
			\begin{tabular}{p{5.5cm}ll}
				\toprule
				Componente & Estatuto & Referência \\
				\midrule
				Soma constante da fita (colunas) &
				Incondicional &
				Prop.~\ref{prop:soma} \\
				Soma das diagonais da fita (par alvo) &
				Incondicional &
				Prop.~\ref{prop:diagonal-alvo} \\
				Bijeção $\Phi$ e soma constante &
				Incondicional &
				Prop.~\ref{prop:phi} \\
				$W_i$ fixa com 4 eixos de simetria &
				Incondicional &
				Def.~\ref{def:Wi} \\
				Soma constante nos eixos de $W_i$ &
				Incondicional &
				Prop.~\ref{prop:soma-Wi} \\
				Par deslocado garantido (Partes~1--2) &
				Incondicional &
				Prop.~\ref{prop:par-deslocado} \\
				Órbita de $\sigma$ / scanner com pivôs &
				Incondicional &
				Prop.~\ref{prop:orbita}; Obs.~\ref{rem:scanner-pivos} \\
				Promoção por deslocamento (Partes~1--2) &
				Incondicional &
				Prop.~\ref{prop:promocao} \\
				Gabarito como referência canônica &
				Incondicional &
				Def.~\ref{def:gabarito} \\
				HI (caso base) &
				Verificável &
				Green--Tao~\cite{green-tao} garante existência \\
				HR$^-$ (par de primos num eixo de $W_i$) &
				Condicional &
				Lacuna central (mínima para o motor) \\
				HR$^+$ (tripla totalmente prima) &
				Condicional &
				Versão forte; implica HR$^-$ \\
				Teorema~\ref{thm:motor} &
				Condicional a HR$^-$ & --- \\
				\bottomrule
			\end{tabular}
		\end{center}
	\end{remark}
	
	\begin{remark}[Cadeia de Markov cíclica]
		O motor com a permutação $\sigma$ tem a estrutura de uma
		\textbf{cadeia determinista cíclica}: o estado é o par
		$(2M,\,\mathcal{E}(t))$ e a transição $\sigma$ é determinística
		na álgebra.  HR$^-$ é equivalente à \textbf{recorrência} da
		cadeia~---~a cadeia não tem estados absorventes, i.e., nenhum
		par $2M$ é evitado para sempre.  Esta conexão com teoria
		ergódica~\cite{grade} será explorada no artigo seguinte.
	\end{remark}
	
	%====================================================================
	\section{O Gabarito~---~$W_i$ Canônica como Referência}
	\label{sec:gabarito}
	%====================================================================
	
	\subsection{O que é o Gabarito}
	
	\begin{definition}[Gabarito]
		\label{def:gabarito}
		O \textbf{gabarito} é uma instância canônica de $W_i$:
		a primeira janela construída pelo scanner na configuração ordenada
		(canônica) da grade~---~isto é, com os ímpares dispostos segundo
		o percurso zigzag $\sigma_G$ (Definição~\ref{def:zigzag-grade})~---~
		que contém pelo menos um par primo nos seus eixos.
		
		Formalmente, seja $W_i^{\text{can}}$ a janela produzida pelo
		scanner na grade canônica que satisfaz: existe um eixo
		$(E_1,E_2)\in\{D_1,D_2,C_v,L_c\}$ com $E_1,E_2\in\mathbb{P}$.
		O gabarito é um objecto fixo~---~referência geométrica~---~e não
		um estado dinâmico.
	\end{definition}
	
	\begin{remark}[Analogia com acelerador de partículas]
		Se o motor fosse um acelerador de partículas, o gabarito seria a
		placa de detecção com os padrões de colisão esperados
		pré-impressos.  Quando os dados dinâmicos se alinham com o
		gabarito, o sensor dispara HR$^-$ ou HR$^+$.
	\end{remark}
	
	\subsection{Por que o Gabarito resolve o problema de estudo}
	
	\begin{remark}[Isolamento do fenômeno]
		Quando $W_i$ é estudada directamente na grade com dados
		aleatórios, os eventos variam conforme a ordenação avança,
		tornando difícil isolar o fenômeno de interesse.  O gabarito
		resolve isto: mantém a estrutura das diagonais fixa como
		referência.
		
		\medskip\noindent\textbf{Estratégia experimental:}
		\begin{enumerate}
			\item Construir o gabarito~---~$W_i^{\text{can}}$ com
			pares obtidos pelo scanner na grade ordenada, com pelo menos
			um par primo.
			\item Iniciar o sistema em entropia máxima (caos puro).
			\item Aplicar ordenação gradual.
			\item A cada passo, sobrepor a configuração dinâmica ao
			gabarito: quando as posições dos eixos coincidem com primos,
			o sensor dispara HR$^-$ ou HR$^+$.
		\end{enumerate}
	\end{remark}
	
	\begin{remark}[Gabarito como referência incondicional]
		A construção do gabarito é inteiramente incondicional:
		a existência dos quatro eixos de $W_i^{\text{can}}$ e a soma
		$E_1+E_2=2M^+$ em cada eixo são garantidas pela
		Proposição~\ref{prop:soma-Wi}.  A primalidade dos extremos~---~
		exigida para que o gabarito exista como referência com pelo menos
		um par primo~---~é a única condição aritmética, verificável
		computacionalmente para qualquer par alvo $2M^+$ específico.
	\end{remark}
	
	%====================================================================
	\section{Protocolo de Simulação~---~Transição Caos $\to$ Ordem}
	\label{sec:simulacao}
	%====================================================================
	
	\subsection{Motivação}
	
	O experimento de transição de fase serve a dois propósitos:
	(1)~demonstrar computacionalmente que HR$^-$ e HR$^+$ são activadas
	durante a ordenação; (2)~medir o \textbf{Tempo de Primeira Passagem}
	(\emph{first-passage time})~---~o número de \emph{swaps} até que a
	primeira janela active o sensor.  Este protocolo é implementado no
	notebook \texttt{motor\_heranca.ipynb}.
	
	\subsection{Protocolo}
	
	\begin{enumerate}
		\item Escolher um par base $2M^-$ e construir a fita-dobra
		$\mathcal{F}_N^1$.
		\item Executar o scanner na grade canônica e construir o gabarito
		$W_i^{\text{can}}$ (referência fixa;
		Definição~\ref{def:gabarito}).
		\item Embaralhar a lista de ímpares: entropia máxima (caos puro).
		\item Verificar no estado inicial: o caos puro já activa o sensor
		por acaso?
		\item Aplicar ordenação total gradual (\emph{Bubble Sort}~---~
		granularidade máxima).
		\item A cada \emph{swap}: executar o scanner e comparar com o
		gabarito.  Registar activação de HR$^-$ e HR$^+$.
		\item Registar: Tempo de Primeira Passagem e percentual do
		percurso total até à activação.
	\end{enumerate}
	
	\subsection{Métricas e Observações Empíricas}
	
	\begin{center}
		\begin{tabular}{p{4.5cm}p{9cm}}
			\toprule
			Métrica & Descrição\\
			\midrule
			Total de \emph{swaps} &
			Número de trocas até à ordenação canônica completa.\\[4pt]
			Passo HR$^-$ (\emph{swap}) &
			\emph{Swap} em que HR$^-$ é activada pela 1.ª vez~---~
			Tempo de Primeira Passagem.\\[4pt]
			\% do percurso &
			Razão passo HR$^-$ / total de \emph{swaps}: quão cedo
			no processo o evento ocorre.\\[4pt]
			HR$^+$ activada &
			Se existe $W_i$ com tripla prima completa
			$(E_1,\,\text{Pivô},\,E_2)$ antes do fim.\\[4pt]
			N.º de janelas $W_i$ &
			Janelas geradas pelo scanner
			$\approx\lceil(N-1)/4\rceil$.\\
			\bottomrule
		\end{tabular}
	\end{center}
	
	\medskip\noindent\textbf{Observações empíricas (testes iniciais,
		$N$ pequeno):}
	\begin{itemize}
		\item HR$^-$ é frequentemente activada já no \emph{swap}~$0$~---~
		no estado de caos puro, antes de qualquer ordenação.  A estrutura
		geométrica de $W_i$ captura pares primos mesmo com dados
		completamente embaralhados.
		\item Quando HR$^-$ não está activa no caos, é activada muito
		cedo no percurso (geralmente abaixo de $10\%$ do total de
		\emph{swaps}), consistente com a Lei de Saturação Geométrica.
		\item HR$^+$ (tripla prima completa) é activada em alguns casos,
		mas não em todos~---~esperado, pois é condição estritamente mais
		restritiva que HR$^-$.
		\item Os testes actuais cobrem $N$ pequeno.  A verificação para
		$N$ grande~---~com múltiplas janelas $W_i$ e análise do
		crescimento do Tempo de Primeira Passagem em função de $N$~---~
		está planeada como próxima etapa computacional.
	\end{itemize}
	
	%====================================================================
	\section{Exemplos Numéricos}
	\label{sec:exemplos}
	%====================================================================
	
	\begin{remark}[Condição de acoplamento nos exemplos]
		O acoplamento $\Phi$ exige $3C=2N$, logo $3\mid N$.  O menor
		exemplo não trivial é $C=8$, $N=12$
		($3\times8=2\times12=24$).  Exemplos com $N$ não divisível por
		$3$ (e.g.\ $N=11$) descrevem a fita isoladamente, sem acoplamento
		a uma grade inteira.
	\end{remark}
	
	\subsection{Fita isolada: $N=11$, $2M^+=44$, $2M^-=42$}
	
	\[
	\begin{array}{c|ccccccccccc}
		k    & 1 & 2 & 3 & 4 & 5 & 6 & 7 & 8 & 9 & 10 & 11\\
		\hline
		a_k  & 1 & 3 & 5 & 7 & 9 &11 &13 &15 &17 &19  &21 \\
		b_k^0& 43&41 &39 &37 &35 &33 &31 &29 &27 &25  &23 \\
		b_k^1& 41&39 &37 &35 &33 &31 &29 &27 &25 &23  &21 \\
	\end{array}
	\]
	
	\begin{itemize}
		\item \textbf{Álgebra (incondicional):} toda coluna soma $44$
		em modo $\varepsilon=0$ e $42$ em modo $\varepsilon=1$,
		independentemente de primalidade.
		\item \textbf{Par base (aritmético):}
		$\mathcal{F}_{11}^0$ tem par primo $(13,31)$ na coluna
		$k=7$ ($13+31=44$~\checkmark);
		$\mathcal{F}_{11}^1$ tem par primo $(13,29)$ na coluna
		$k=7$ ($13+29=42$~\checkmark).
	\end{itemize}
	
	\subsection{Acoplamento completo: $C=8$, $N=12$, $2M^+=48$,
		$2M^-=46$}
	
	$3\times8=2\times12=24$ células. \checkmark
	
	\medskip\noindent\textbf{Grade $G_{3,8}$:}
	\[
	\begin{array}{c|cccccccc}
		& c=1&2&3&4&5&6&7&8\\
		\hline
		\text{Linha~1 }(\rightarrow) & 1&3&5&7&9&11&13&15\\
		\text{Linha~2 }(\leftarrow)  &31&29&27&25&23&21&19&17\\
		\text{Linha~3 }(\rightarrow) &33&35&37&39&41&43&45&47\\
	\end{array}
	\]
	
	\medskip\noindent\textbf{Fita $\mathcal{F}_{12}$:}
	\[
	\begin{array}{c|cccccccccccc}
		k & 1&2&3&4&5&6&7&8&9&10&11&12\\
		\hline
		a_k   &1&3&5&7&9&11&13&15&17&19&21&23\\
		b_k^0 &47&45&43&41&39&37&35&33&31&29&27&25\\
		b_k^1 &45&43&41&39&37&35&33&31&29&27&25&23\\
	\end{array}
	\]
	
	\medskip\noindent\textbf{Acoplamento $\Phi$~---~camada algébrica
		(incondicional):}
	\[
	\begin{array}{c|cc|c|c}
		k & \Phi(k)_1 & \Phi(k)_2 & \text{soma} & \text{células em }G_{3,8}\\
		\hline
		1  & 1  & 47 & 48 & f(1,1)+f(3,8)\\
		2  & 3  & 45 & 48 & f(1,2)+f(3,7)\\
		3  & 5  & 43 & 48 & f(1,3)+f(3,6)\\
		4  & 7  & 41 & 48 & f(1,4)+f(3,5)\\
		5  & 9  & 39 & 48 & f(1,5)+f(3,4)\\
		6  & 11 & 37 & 48 & f(1,6)+f(3,3)\\
		7  & 13 & 35 & 48 & f(1,7)+f(3,2)\\
		8  & 15 & 33 & 48 & f(1,8)+f(3,1)\\
		9  & 17 & 31 & 48 & f(2,8)+f(2,1)\\
		10 & 19 & 29 & 48 & f(2,7)+f(2,2)\\
		11 & 21 & 27 & 48 & f(2,6)+f(2,3)\\
		12 & 23 & 25 & 48 & f(2,5)+f(2,4)\\
	\end{array}
	\]
	Toda coluna soma $48=4N=2M^+$, independentemente de primalidade. \checkmark
	
	\medskip\noindent\textbf{Motor em acção~---~camada aritmética
		(condicional a HR$^-$):}
	\begin{itemize}
		\item $\mathcal{F}_{12}^1$ ($2M^-=46$): par base primo $(5,41)$
		na coluna $k=3$ ($5+41=46$~\checkmark).
		O scanner regista o par $(7,41)$ em $t=1$ num eixo de $W_i$
		($C=8$): ambos primos e $7+41=48=2M^+$~\checkmark~---~evento
		restritivo HR$^-$.  O par $t=2$ dá $(11,37)$:
		$11+37=48$~\checkmark, também evento restritivo.
		\item \textbf{Promoção (incondicional):} o par $(7,41)$ é
		promovido a par base para $2M^+=48$ pelo mecanismo de
		Proposição~\ref{prop:promocao}~---~Goldbach para
		$48$~\checkmark.
		\item Par prometido para $2M^+=50$: determinado pelo próximo
		evento restritivo em $G_{3,C'}$ após deslocamento mínimo
		(requer HR$^-$ na nova configuração).
	\end{itemize}
	
	%====================================================================
	\section{Conclusão}
	\label{sec:conclusao}
	%====================================================================
	
	Este artigo constitui a \textbf{fundação algébrica} do Motor de
	Herança Estrutural na série sobre a Conjectura de Goldbach.
	
	\medskip
	\noindent\textbf{Camada algébrica~---~resultados incondicionais.}
	Demonstrámos por álgebra pura, sem qualquer hipótese de primalidade:
	\begin{enumerate}
		\item A \textbf{soma constante das colunas}
		$a_k+b_k=2M(\varepsilon)$ em toda coluna da fita~---~propriedade
		posição-cega e primalidade-cega (Proposição~\ref{prop:soma}).
		\item A \textbf{soma constante das diagonais}
		$a_{k+1}+b_k^{(1)}=2M^+$ (par alvo): as diagonais da fita do
		par base somam exactamente o par alvo, incondicionalmente
		(Proposição~\ref{prop:diagonal-alvo}).  Esta distinção~---~
		colunas $=$ par base, diagonais $=$ par alvo~---~é a distinção
		central do motor.
		\item A \textbf{bijeção} $\Phi$ entre $G_{3,C}$ e
		$\mathcal{F}_N$, com soma constante $\Phi(k)_1+\Phi(k)_2
		=4N$ para todo~$k$ (Proposição~\ref{prop:phi}).
		\item A \textbf{estrutura de $W_i$} como matriz $3\times3$ com
		quatro eixos de simetria ($D_1$, $D_2$, $C_v$, $L_c$), cada um
		armazenando um par diagonal com soma $2M^+$~---~garantida
		incondicionalmente pela construção da fita
		(Definição~\ref{def:Wi}; Proposição~\ref{prop:soma-Wi}).
		\item O \textbf{par deslocado garantido}: os extremos de qualquer
		eixo de $W_i$ somam $2M^+$ incondicionalmente (via acoplamento
		$\Phi$), e o próximo deslocamento mínimo propaga essa soma para
		$2M^++2$~---~quer os extremos sejam ou não primos
		(Proposição~\ref{prop:par-deslocado}).
		\item A \textbf{órbita completa} de $\sigma$ (equivalente ao
		scanner com pivôs): os $N-1$ pares diagonais são cobertos em
		bijecção, com soma constante $2M^+$ em todo passo
		(Proposição~\ref{prop:orbita}; Observação~\ref{rem:scanner-pivos}).
		\item A \textbf{estrutura direccional} de $\sigma$ e a
		\textbf{identidade dos pares diagonais do scanner}: L1 e L3
		movem-se em sentidos opostos ($e_1\!\uparrow$,
		$e_3\!\downarrow$), os extremos de $W_i$ apontam
		exactamente os pares diagonais $(a_{k+1},b_k^{(1)})$ com soma
		$2M^+$, e o passo unitário $\delta=1$ é a única escolha
		que mantém simultaneamente a cobertura bijectiva e a
		sincronização $N\to N{+}1$
		(Observação~\ref{rem:sigma-direcoes};
		Proposição~\ref{prop:pares-diagonais}).
		\item O \textbf{mecanismo de promoção}: qualquer deslocamento
		$+2$ na fita propaga-se para $G_{3,C}$ como operação de
		pilha global, promovendo qualquer par simétrico~---~primo
		ou não~---~ao alinhamento com a fita seguinte
		(Proposição~\ref{prop:promocao}).
		\item O \textbf{gabarito} $W_i^{\text{can}}$ como referência
		canônica fixa: a janela obtida na configuração ordenada da grade,
		contra a qual os estados dinâmicos são comparados durante a
		simulação (Definição~\ref{def:gabarito}).
	\end{enumerate}
	
	\medskip
	\noindent\textbf{Camada aritmética~---~resultados condicionais.}
	Sob as Hipóteses Restritivas HR$^-$/HR$^+$, que postulam a
	existência de pelo menos um passo $t$ da órbita de $\sigma$ em que
	os extremos de $W_i$ são ambos primos (HR$^-$) ou a tripla
	completa é totalmente prima (HR$^+$):
	\begin{enumerate}
		\item O \textbf{par prometido} existe
		(Definição~\ref{def:prometido}): os seus extremos, cujo
		alinhamento é garantido algebricamente, são de primos por
		hipótese.
		\item O \textbf{Teorema do Motor}
		(Teorema~\ref{thm:motor}) estabelece que a cadeia
		$2M_0\to2M_0+2\to\cdots$ cobre todos os pares $\geq 2M_0$
		como somas de dois primos.
	\end{enumerate}
	
	\medskip
	\noindent\textbf{Lacuna remanescente.}
	A única hipótese não provada é a primalidade simultânea dos
	extremos $(e_1(t),e_3(t))$ em algum passo~$t$~---~i.e., HR$^-$.
	A existência geométrica do par e a sua cobertura pela órbita de
	$\sigma$ são incondicionais; a primalidade é o único ingrediente
	aritmético.  HR$^-$ (e a versão mais forte HR$^+$) serão o objecto
	do artigo seguinte, que as atacará usando o peso oracle-free
	$\Omega_N^*$~\cite{peso} e a desigualdade estrutural
	$\mu_{\text{comp}}>\mu_{\text{global}}>\mu_{\text{primo}}$~\cite{problemaB},
	com o acoplamento $\Phi$ e a órbita de $\sigma$ disponíveis como
	ferramentas algébricas.
	
	%====================================================================
	\begin{thebibliography}{9}
		%====================================================================
		
		\bibitem{grade}
		T.~Bandeira,
		\textit{Uma Abordagem Unificada para a Conjectura de Goldbach:
			Estrutura Combinatória, Saturação Geométrica e o Elo entre Trios
			e Pares Primos}, preprint, maio de 2026.
		
		\bibitem{invariante}
		T.~Bandeira,
		\textit{Uma propriedade de soma constante na grade $3\times N$
			e as conjecturas de Goldbach}, preprint, maio de 2026.
		
		\bibitem{peso}
		T.~Bandeira,
		\textit{O Invariante Âncora da Grade $G_{3,N}$ e um Estimador
			Geométrico Oracle-Free para a Conjectura de Goldbach},
		preprint, maio de 2026.
		
		\bibitem{problemaB}
		T.~Bandeira,
		\textit{Problema~B do Crivo Geométrico: Demonstração da
			Desigualdade Estrutural}, preprint, maio de 2026.
		
		\bibitem{motor}
		T.~Bandeira,
		\textit{Motor de Herança Estrutural em Grades $3\times N$:
			Um Sistema Dinâmico de Cobertura Condicional para Pares de
			Goldbach}, preprint, maio de 2026.
		
		\bibitem{green-tao}
		B.~Green e T.~Tao,
		\textit{The primes contain arbitrarily long arithmetic
			progressions},
		\textit{Annals of Mathematics}, vol.~167, pp.~481--547, 2008.
		
	\end{thebibliography}
	
\end{document}